\documentclass[12pt, twocolumn]{article}

%-------Packages---------
\usepackage{newpxtext,newpxmath} % Palatino-like text & math (modern)
\let\openbox\relax
\usepackage{amsthm}
\usepackage{bboldx}
\usepackage{mathtools}
\usepackage{tikz}
\usetikzlibrary{calc,intersections}
\usepackage{tikz-cd}
\usepackage[all,arc]{xy}
\usepackage{graphicx}

\usepackage{enumerate}
\usepackage[dvipsnames]{xcolor}
\usepackage{float}
\usepackage[left=0.4in,top=0.6in,right=0.4in,bottom=0.4in]{geometry}
\usepackage[nocheck]{fancyhdr}
\usepackage{multicol}
\usepackage{setspace}

\usepackage{dsfont}
\usepackage{mathrsfs}

\usepackage{tcolorbox}
\tcbuselibrary{theorems,skins,breakable}

\usepackage{xparse}
\usepackage{import}
\usepackage[utf8]{inputenc}

\usepackage{hyperref}
\usepackage{cleveref}

\documentclass[12pt, letterpaper]{article}
\usepackage{xcolor, tcolorbox, amsthm, amsmath}
\tcbuselibrary{theorems,skins,breakable}
% Theorem System
% The following boxes are provided:
%   Definition:     \defn 
%   Theorem:        \thm 
%   Lemma:          \lem
%   Corollary:      \cor
%   Proposition:    \prop   
%   Claim:          \clm
%   Proof:          \pf
%   Example:        \ex
%   Remark:         \rmk (sentence), \rmkb (block)
% Suffix
%   r:              Allow Theorem/Definition to be referenced, e.g. thmr
%   p:              Add a short proof block for Lemma, Corollary, Proposition or Claim, e.g. lemp
%                   For theorems, use \pf for proof blocks

% ======= Real examples : 

% \defn{Definition Name}{
%     A defintion.
% }

% \thmr{Theorem Name}{mybigthm}{
%     A theorem.
% }

% \lem{Lemma Name}{
%     A lemma.
% }


% \cor{
%     A corollary.
% }

% \prop{
%     A proposition.
% }

% \clmp{}{
%     A claim.
% }{
%     A reference to Theorem~\ref{thm:mybigthm}
% }

% \pf{
%    proof.

% \ex{
%     some examples. 
% }

% \rmk{
%     Some remark.
% }

% \rmkb{
%     Some more remark.
% }
% Define custom colors
\definecolor{thmscol}{HTML}{F4565B} %For Theorems
\definecolor{lemscol}{HTML}{FE844C} %For Lemmes
\definecolor{prpscol}{HTML}{00A7EF} %For proposition
\definecolor{corscol}{HTML}{23DAFF} %For corrolaries
\definecolor{rmkscol}{HTML}{6DEEC5} %For remarks

\definecolor{defscol}{HTML}{18C991} %For definitions

\definecolor{exmscol}{HTML}{7DB800} %For examples

\definecolor{clmscol}{HTML}{165c58} %For claims


% ============================
% Definition
% ============================
\newtcbtheorem[number within=section]{definition}{Definition}
{
    enhanced,
    frame hidden,
    titlerule=0mm,
    toptitle=1mm,
    bottomtitle=1mm,
    fonttitle=\bfseries\large,
    coltitle=black,
    colbacktitle=defscol!70!white,
    colback=defscol!35!white,
}{defn}

\NewDocumentCommand{\defn}{m+m}{
    \begin{definition}{#1}{}
        #2
    \end{definition}
}

\NewDocumentCommand{\defnr}{mm+m}{
    \begin{definition}{#1}{#2}
        #3
    \end{definition}
}

% ============================
% Theorem
% ============================

\newtcbtheorem[use counter from=definition]{theorem}{Theorem}
{
    enhanced,
    frame hidden,
    titlerule=0mm,
    toptitle=1mm,
    bottomtitle=1mm,
    fonttitle=\bfseries\large,
    coltitle=black,
    colbacktitle=thmscol!80!white,
    colback=thmscol!38!white,
}{thm}

\NewDocumentCommand{\thm}{m+m}{
    \begin{theorem}{#1}{}
        #2
    \end{theorem}
}

\NewDocumentCommand{\thmr}{mm+m}{
    \begin{theorem}{#1}{#2}
        #3
    \end{theorem}
}

\newenvironment{thmpf}{
	{\noindent{\it \textbf{Proof.}}}
	\tcolorbox[blanker,breakable,left=5mm,parbox=false,
    before upper={\parindent15pt},
    after skip=10pt,
	borderline west={1mm}{0pt}{thmscol!80!white}]
	\setlength{\parindent}{0pt}
}{
    \textcolor{thmscol!80!white}{\hbox{}\nobreak\hfill$\blacksquare$} 
    \endtcolorbox
}

\NewDocumentCommand{\thmp}{m+m+m}{
    \begin{theorem}{#1}{}
        #2
    \end{theorem}

    \begin{thmpf}
        #3
    \end{thmpf}
}

% ============================
% Lemma
% ============================

\newtcbtheorem[use counter from=definition]{lemma}{Lemma}
{
    enhanced,
    frame hidden,
    titlerule=0mm,
    toptitle=1mm,
    bottomtitle=1mm,
    fonttitle=\bfseries\large,
    coltitle=black,
    colbacktitle=lemscol!80!white,
    colback=lemscol!40!white,
}{lem}

\NewDocumentCommand{\lem}{m+m}{
    \begin{lemma}{#1}{}
        #2
    \end{lemma}
}

\newenvironment{lempf}{
	{\noindent{\it \textbf{Proof.}}}
	\tcolorbox[blanker,breakable,left=5mm,parbox=false,
    before upper={\parindent15pt},
    after skip=10pt,
	borderline west={1mm}{0pt}{lemscol!50!white}]
	\setlength{\parindent}{0pt}
}{
    \textcolor{lemscol!50!white}{\hbox{}\nobreak\hfill$\blacksquare$} 
    \endtcolorbox
}

\NewDocumentCommand{\lemp}{m+m+m}{
    \begin{lemma}{#1}{}
        #2
    \end{lemma}

    \begin{lempf}
        #3
    \end{lempf}
}

% ============================
% Corollary
% ============================

\newtcbtheorem[use counter from=definition]{corollary}{Corollary}
{
    enhanced,
    frame hidden,
    titlerule=0mm,
    toptitle=1mm,
    bottomtitle=1mm,
    fonttitle=\bfseries\large,
    coltitle=black,
    colbacktitle=corscol!50!white,
    colback=corscol!20!white,
}{cor}

\NewDocumentCommand{\cor}{+m}{
    \begin{corollary}{}{}
        #1
    \end{corollary}
}

\newenvironment{corpf}{
	{\noindent{\it \textbf{Proof.}}}
	\tcolorbox[blanker,breakable,left=5mm,parbox=false,
    before upper={\parindent15pt},
    after skip=10pt,
	borderline west={1mm}{0pt}{corscol!50!white}]
	\setlength{\parindent}{0pt}
}{
    \textcolor{corscol!50!white}{\hbox{}\nobreak\hfill$\blacksquare$} 
    \endtcolorbox
}

\NewDocumentCommand{\corp}{m+m+m}{
    \begin{corollary}{#1}{}
        #2
    \end{corollary}

    \begin{corpf}
        #3
    \end{corpf}
}

% ============================
% Proposition
% ============================

\newtcbtheorem[use counter from=definition]{proposition}{Proposition}
{
    enhanced,
    frame hidden,
    titlerule=0mm,
    toptitle=1mm,
    bottomtitle=1mm,
    fonttitle=\bfseries\large,
    coltitle=black,
    colbacktitle=prpscol!70!white,
    colback=prpscol!20!white,
}{prop}

\NewDocumentCommand{\prop}{+m}{
    \begin{proposition}{}{}
        #1
    \end{proposition}
}

\newenvironment{proppf}{
	{\noindent{\it \textbf{Proof.}}}
	\tcolorbox[blanker,breakable,left=5mm,parbox=false,
    before upper={\parindent15pt},
    after skip=10pt,
	borderline west={1mm}{0pt}{prpscol!50!white}]
	\setlength{\parindent}{0pt}
}{
    \textcolor{prpscol!50!white}{\hbox{}\nobreak\hfill$\blacksquare$} 
    \endtcolorbox
}



\NewDocumentCommand{\propp}{+m+m}{
    \begin{proposition}{}{}
        #1
    \end{proposition}

    \begin{proppf}
        #2
    \end{proppf}
}

% ============================
% Claim
% ============================

\newtcbtheorem[use counter from=definition]{claim}{Claim}
{
    enhanced,
    frame hidden,
    titlerule=0mm,
    toptitle=1mm,
    bottomtitle=1mm,
    fonttitle=\bfseries\large,
    coltitle=black,
    colbacktitle=clmscol!50!white,
    colback=clmscol!20!white,
}{clm}

\NewDocumentCommand{\clm}{m+m}{
    \begin{claim*}{#1}{}
        #2
    \end{claim*}
}

\newenvironment{clmpf}{
	{\noindent{\it \textbf{Proof.}}
	\setlength{\parindent}{0pt}
	}
	\tcolorbox[blanker,breakable,left=5mm,parbox=false,
    before upper={\parindent15pt},
    after skip=10pt,
	borderline west={1mm}{0pt}{clmscol!50!white}]
}{
    \textcolor{clmscol!50!white}{\hbox{}\nobreak\hfill$\blacksquare$} 
    \endtcolorbox
}

\NewDocumentCommand{\clmp}{m+m+m}{
    \begin{claim*}{#1}{}
        #2
    \end{claim*}

    \begin{clmpf}
        #3
    \end{clmpf}
}


% ============================
% Proof
% ============================

\NewDocumentCommand{\pf}{+m}{
    \begin{proof}
        [\noindent\textbf{Proof.}]
        #1
    \end{proof}
}

% ============================
% Example
% ============================

\newtcbtheorem[use counter from=definition]{example}{Example}
{
    blanker,
    breakable,
    left=5mm,
    parbox=false,
	borderline west={1mm}{0pt}{exmscol!50!white},
    titlerule=0mm,
    toptitle=1mm,
    bottomtitle=1mm,
    fonttitle={\bfseries\itshape},
    after skip=10pt,
    top=1mm,
    coltitle=black
}{exm}

\NewDocumentCommand{\exm}{+m+m}{
    \begin{example*}{#1}{}
        #2
        \setlength{\parindent}{0pt}
        \textcolor{exmscol!50!white}{\hbox{}\nobreak\hfill$\blacksquare$}
    \end{example*}
}


% ============================
% Remark
% ============================


\NewDocumentCommand{\rmk}{+m}{
    {\it \color{rmkscol!100!white}#1}
}

\newenvironment{remark}{
    \par
    \vspace{5pt}
    \begin{minipage}{\textwidth}
        {\par\noindent{\textbf{Remark.}}}
        \tcolorbox[blanker,breakable,left=5mm,
        before skip=10pt,after skip=10pt,
        borderline west={1mm}{0pt}{rmkscol!80!white}]
}{
        \endtcolorbox
    \end{minipage}
    \vspace{5pt}
}

\NewDocumentCommand{\rmkb}{+m}{
    \begin{remark}
        #1
    \end{remark}
}

\dump


\definecolor{thmscol}{HTML}{FFFFFF} %For Theorems
\definecolor{defscol}{HTML}{FFFFFF} %For definitions
\definecolor{exmscol}{HTML}{FFFFFF} %For examples

% Define a new command for problems
\renewcommand{\headrule}{\color{black}\hrulefill}
\newcommand{\C}{\mathbb{C}}
\newcommand{\R}{\mathbb{R}}
\newcommand{\Q}{\mathbb{Q}}
\newcommand{\Z}{\mathbb{Z}}
\newcommand{\N}{\mathbb{N}}
\renewcommand{\P}{\mathbb{P}}
\newcommand{\F}{\mathbb{F}}
\newcommand{\contr}{\Rightarrow \Leftarrow}
\newcommand{\s}{\(\,\)}
\renewcommand{\t}[1]{\text{#1}}

\newcommand{\skcgen}{{\sf Gen}} 
\newcommand{\skcenc}{{\sf Enc}}
\newcommand{\skcdec}{{\sf Dec}}
\newcommand{\mac}{{\sf MAC}}
\newcommand{\ver}{{\sf Vrfy}}

% Meta data
\setstretch{1.2}

\begin{document}
\pagecolor{white}
\setlength{\parindent}{0pt}
\pagestyle{fancy}
\fancyhf{}
\fancyhead[LOE]{\slshape{\textbf{Vincent Ferrara}}}
\fancyhead[ROE]{Midterm --- \thepage}
\fancyhead[C]{EECS 475}

\section{Perfect Secrecy}
\vspace{-0.2in}
\defn{Correctness}{
    $\forall k \in \mathrm{Im}(\skcgen), \forall m \in \mathcal{M},$
    \vspace{-0.2in}
    \[\skcdec(k,\skcenc(k,m))=m\]
}
\vspace{-0.2in}
\defn{Shannon \\Secrecy}{
    An encryption scheme $(\skcgen, \skcenc, \skcdec)$ with message space $\mathcal{M}$
    is said to be Shannon Secret if for every probability distribution over
    $\mathcal{M}$, $\forall m \in \mathcal{M}$, and $\forall c \in \mathcal{C}$
    for which $\P[c \in \mathcal{C}] > 0$,
    \vspace{-0.2in}
    \[\P[M = m \mid C =c ] = \P[M=m]\]
}
\vspace{-0.2in}
\defn{Perfect Secrecy}{
    An encryption scheme $(\skcgen, \skcenc, \skcdec)$ with message space $\mathcal{M}$
    is said to be Perfectly Secret if $\forall m_0, m_1 \in \mathcal{M}, \forall c \in \mathcal{C}$
    \vspace{-0.2in}
    \[\P[C = c \mid M=m_0] = \P[C = c \mid M=m_1]\]
}
\vspace{-0.2in}
\thm{}{
    Perfect Secrecy is equivalent to Shannon Secrecy
}
\vspace{-0.2in}
\thm{}{
    If an encryption scheme $(\skcgen, \skcenc, \skcdec)$ is perfectly secret
    with message space $\mathcal{M}$ and key space $\mathcal{K}$, then $|\mathcal{K}| \geq |\mathcal{M}|$
}
\vspace{-0.2in}
\defn{One-Time Pad}{
    $\mathcal{M}=\mathcal{C}=\mathcal{K}=\{0,1\}^\ell, \t{ where } \ell \geq 1$\\
    $\skcgen : k \overset{\$}{\leftarrow} \mathcal{K}$, 
    
    $\skcenc(k,m) : c = m \oplus k$, 
    
    $\skcdec(k,m) : c \oplus k$
}
\vspace{-0.4in}
\subsection{Flaws with perfect secrecy}
\begin{itemize}
    \item Can't use same key twice.
    \item Key must be chosen uniformly at random and be at least as long as the message.
\end{itemize}

\section{Negligible Functions}
\vspace{-0.2in}
\defn{Negligible Function}{
    A non-negative function $\epsilon$ is negligible if for all $c \geq 0$ there exists
    a $N \in \N$ s.t. for all $n \geq N$
    \vspace{-0.2in}
    \[\epsilon(n) \leq 1/n^c\]
}
\vspace{-0.2in}
\thm{}{
    For $f,g$ polynomially bounded functions and negligible functions $\epsilon$ and $\phi$.
    \vspace{-0.15in}
    \begin{itemize}
        \item $f \cdot \epsilon = \eta$
        \vspace{-0.2in}
        \item $f+g=h$, $f \cdot g = h'$
        \vspace{-0.2in}
        \item $\epsilon + \phi = \eta'$, $\epsilon \cdot \phi = \eta''$
    \end{itemize}
    \vspace{-0.2in}
    For $\eta,\eta',\eta''$ negligible functions and $h,h'$ polynomially bounded functions.
}
\vspace{-0.45in}
\section{Computationally Secure}
\vspace{-0.2in}
\defn{Computationally\\ Indistinguishable}{
    \small Given ensembles $\{X_n^0\}_{n \in \N}$ and $\{X_n^1\}_{n \in \N}$
    where $X_n$ are distributions over $\{0,1\}^{\ell(n)}$ where $\ell$ is some polynomial,
    we say that $\{X_n^0\}_{n \in \N}$ and $\{X_n^1\}_{n \in \N}$ are computationally
    Indistinguishable iff $\forall $ PPT $\mathcal{A}$, there exists a $\mathrm{negl}(\cdot)$ s.t.
    for all $n \in \N$,

    (Predictor):
    \vspace{-0.2in}
    \[\P[b \leftarrow \{0,1\}, x \leftarrow X^b_n \mid \mathcal{A}(x)=b] \leq 1/2 + \mathrm{negl}(n)\]
    
    \vspace{-0.2in}
    (Decider):
    \vspace{-0.2in}
    \[|\P[x \leftarrow X_n^0 \mid \mathcal{A}(x)=1] - \P[x \leftarrow X_n^1 \mid \mathcal{A}(x)=1]|\] 
    \vspace{-0.4in}
    \[\leq \mathrm{negl}(n)\]
}
\vspace{-0.2in}
\thm{}{
    Computationally Indistinguishable ensembles are closed under efficient operations and transitive up to a polynomial amount.
}
\vspace{-0.2in}
\defn{EAV Security}{
    $\mathsf{EXPT}_{EAV}^{b, \mathcal{A}}(1^n)$: $\mathcal{A}$ sends two messages of the same length $(m_0,m_1)$ then the challenger chooses one of the two messages depending on the parameter $b$
    then returns $c_b$ to the $\mathcal{A}$. Then $\mathcal{A}$ outputs its guess $b'$.
}
\section{Pseudorandom Generators}
\vspace{-0.15in}
\defn{PRG}{
    A determ. algo $G : \{0,1\}^n \to \{0,1\}^{\ell(n)} \in \mathsf{P}$ is a PRG if
    \vspace{-0.2in}
    \begin{itemize}
        \item expansion: $\forall n, \ell(n) > n$.
        \vspace{-0.2in}
        \item pseudoran: $\{s \leftarrow U_n \mid G(s)\}_n \approx U_{\ell(n)}$
    \end{itemize}
}
\vspace{-0.2in}
\defn{PRG One-Time Pad}{
    Same as One-Time pad but instead of $\oplus k$ do $\oplus G(k)$.
    Can't reuse key.
}
\vspace{-0.45in}
\section{Multi-msg Security}
\vspace{-0.2in}
\defn{Multi-EAV Security}{
    $\mathsf{EXPT}^{b,\mathcal{A}}_{\mathrm{Meav}}(1^n)$:
    $\mathcal{A}$ sends two vectors of messages where $|m_{0,i}|=|m_{1,i}|$ for all $i$. 
    The challenger chooses one of the two vectors and returns one back encrypted. $\mathcal{A}$ outputs its guess $b'$.
}
\vspace{-0.2in}
\defn{Multi-CPA Security}{
    $\mathsf{EXPT}^{b,\mathcal{A}}_{\mathrm{LRcpa}}(1^n)$:
    $\mathcal{A}$ can query the oracle $\mathrm{LR}_k^b$ a polynomial times with $m_0^i,m_1^i$ and each time the oracle will
    always return either the left message encryption or the right message. Then $\mathcal{A}$ outputs its guess. 
    
}
\vspace{-0.2in}
\defn{CPA Security}{
    $\mathsf{EXPT}^{b,\mathcal{A}}_{\mathrm{cpa}}(1^n)$: This basically the Multi-CPA exp. 
    First make queries then input actual challenge it sends back then query then guess.
}
\vspace{-0.2in}
\thm{}{
    Multi-CPA $\implies$ Multi-EAV\\
    Multi-CPA $\iff$ CPA
}
\vspace{-0.2in}
\thm{}{
    In order for a scheme to be CPA secure, it must have randomized encryption.
}
\section{Pseudorandom Functions}
\vspace{-0.15in}
\defn{PRF}{
    A determ. algo $F : \mathcal{K} \times \{0,1\}^n \to \{0,1\}^n \in \mathsf{P}$ is a PRF if
    
    $\{k \leftarrow \mathcal{K}, x \leftarrow U_n \mid F_k(x)\} \approx$ 
    \\$\{x \leftarrow U_n \mid F(x)\} \approx U_n$ where $F$ is a uniformly random function.
}
\vspace{-0.2in}
\thm{}{
    PRG $\iff$ PRF
}
\vspace{-0.2in}
\defn{CPA secure with PRF}{
    Let $F$ be PRF. 
    
    $\skcgen(1^n): k \overset{\$}{\leftarrow} \{0,1\}^n$

    $\skcenc(k,m): r \overset{\$}{\leftarrow} \{0,1\}^n, c = <r, F_k(r) \oplus m>$

    $\skcdec(k,<r,c'>): m := F_k(r) \oplus c'$
}
\vspace{-0.45in}
\section{MAC}
\vspace{-0.15in}
\defn{EU MAC}{
    $\mathsf{EXPT}_{EU}^\mathcal{A}(1^n)$
    $\mathcal{A}$ queries then guesses $(m,t)$.
    If $\mathsf{Vrfy}(k,m,t)=1$ and $m$ not queried by $\mathcal{A}$ then output 1.
}
\vspace{-0.2in}
\defn{SEU MAC}{
    $\mathsf{EXPT}_{SEU}^\mathcal{A}(1^n)$
    Same game but succeeds if the guess $(m,t)$ has not been seen before.
}
\vspace{-0.2in}
\thm{}{
    If a MAC with canonical verification EU then it is SEU

    $\mathrm{Vrfy}(k,m,t)=1$ if $\mathrm{MAC}(k,m)=t$ this implies MAC is determinitistic
}
\vspace{-0.2in}
\defn{PRF MAC}{
    Let $F$ be a PRF.
    
    $\mathsf{MAC}(k,m) = F_k(m) = t$

    $\mathsf{Vrfy}(k,m,t) = $ Canonical
}
\defn{MAC for Arb len}{
    Let MAC' be secure for $n$-bit msg.
    
    $\mathsf{MAC}(k,m)$: part $m=m_1||m_2||\dots||m_d$, where $|m_i|=n/4$. Let $r \overset{\$}{\leftarrow} U_{n/4}$,
    then let $t_i = \mathsf{MAC}'(k, r||l||i||m_i)$ then let $t=(r,t_1,\dots,t_d)$.
}
\vspace{-0.45in}
\section{Hash Functions}
\vspace{-0.15in}
\defn{Hash Family}{
    A determ. algo $H: \mathcal{K} \times \{0,1\}^* \to \{0,1\}^n$ s.t. $*=$Arb len.
}
\vspace{-0.2in}
\defn{Collision Resistant}{
    $\mathsf{EXPT}^\mathcal{A}(1^n)$ The challenger sends the key $s$ then the has polynomial time to send back $(x,y)$ s.t. $H^s(x)=H^s(y)$.
    succeeds if it finds a pair.
}
\vspace{-0.2in}
\defn{Hash MAC}{
    Let MAC' be a secure $n$-bit msg.

    $\mathsf{MAC}((k,s),m)=\mathsf{MAC'}(H^s(m))$
}
\[\begin{tikzcd}
	&& {0,3} \\
	& {0,1} & {2,3} \\
	0 & 1 & 2 & 3
	\arrow[from=2-2, to=1-3]
	\arrow[from=2-3, to=1-3]
	\arrow[from=3-1, to=2-2]
	\arrow[from=3-2, to=2-2]
	\arrow[from=3-3, to=2-3]
	\arrow[from=3-4, to=2-3]
\end{tikzcd}\]
Merkle path is the nodes you needed to construct the root.

User only has to store the root digest to verify all of tree. To update a block for example block 1
all the server has to send to you is block 0 and 2,3. Also, same with verifying a block all you need is the block 1, 0, 2,3 and the user stores the root, so they can verify.

To check and make more efficient first make it, so you can lose half the blocks to then lose data then use a probabilistic check to get a probability of $1/2^{-k}$ of passing if they did change one block.

\section{Number Theory}
\vspace{-0.2in}
\thm{Fermat's Little Thm}{
    If $p$ is prime then for $a \in \{1,\dots,p-1\}$ then $a^{p-1} \equiv 1 \mod p$
}
\vspace{-0.2in}
\exm{Pick Subgroups of Prime Order}{
    Let $p=rq+1$ where $p,q$ are primes.

    Then $\mathbb{G}=\{[h^r \mod p] \mid h \in \Z_p^*\}$ is a subgroup of order $q$.
}
\vspace{-0.3in}
\section{Key Exchange}
\vspace{-0.2in}
\defn{Dlog}{
    $\mathsf{Dlog}^{\mathcal{A},\mathcal{G}}(1^n)$
    Challenger sends $(\mathbb{G},q,g,h)$ where $h=g^x$ and $\mathcal{A}$ guesses $x$.
    If correct then return 1.
}
\vspace{-0.2in}
\defn{DDH}{
    $\mathsf{DDH}_b^{\mathcal{A},\mathcal{G}}(1^n)$
    Challenger sends $(\mathcal{G},q,g,g^x,g^y,h_b)$ s.t. $h_0=g^z$ and $h_1=g^{xy}$. Then $\mathcal{A}$ guesses if $h_b$ is random or $g^{xy}$.
}
\vspace{-0.15in}
DDH is harder than Dlog.
\vspace{-0.15in}
\defn{DH Key Exchange}{
    First person samples a prime order multiplicative group $(\mathbb{G}, q,g)$ and samples an $x \in \Z_q$ then they send $(\mathbb{G}, q,g),g^x$ to the second person.
    The second person then samples $y \in \Z_q$ and then sends $g^y$ back then they both can compute $g^{xy}$.
}
\vspace{-0.2in}
\defn{Key Expt}{
    $\mathsf{KEEXPT}_b(1^n)$
    Two parties run the key exchange protocol. Let $\Gamma$ be the transcript of the protocol and $k$ the outputted key at the end.
    If $b=0$ output $\gamma, k$ and else $\gamma, k'$ where $k' \overset{\$}{\leftarrow} \mathcal{K}$.
}
\vspace{-0.45in}
\section{Public-Key Crypt}
\vspace{-0.15in}
\subsection{Syntax}
$\skcgen(1^n) \to (pk,sk)$
$\skcenc(pk,m) \to c$
$\skcdec(sk,c)=m$
\defn{Correctness}{
    Except with negligible probability over the choice of $(pk,sk)$ output by $\skcgen(1^n)$, it must
    be that for any legal msg $m$,
    \vspace{-0.2in}
    \[\skcgen(sk, \skcenc(pk,m))=m\]
}
\vspace{-0.2in}
\defn{EAV}{
    $\mathsf{PKEXPT}_{eav}^{b,\mathcal{A}}(1^n)$
    Challenger sends a public key $pk$ and then $\mathcal{A}$ sends $m_0,m_1$ and challenger picks one to encrypt and send back.
}
\vspace{-0.2in}
\thm{}{
    A public key encryption protocol is EAV-secure if and only if it is CPA-secure.

    This implies it must be random.
}
\vspace{-0.2in}
\defn{El Gamel Enc}{
    Run the  group protocol to get a prime order multiplicative group $(\mathbb{G},q,g)$.

    $\skcgen$: sample $x \leftarrow \Z_q$ and output $pk=(\mathbb{G},q,h=g^x)$ and $sk=(\mathbb{G},q,g,x)$.

    $\skcenc(pk, m)$: choose random $y \leftarrow \Z_q$ and output $(g^y, m h^y)$

    $\skcdec(sk,(c_1=g^y,c_2=mg^{xy}))$: $c_2 \cdot g^{-xy} = m$
}
\vspace{-0.45in}
\section{Signature}
\vspace{-0.2in}
\subsection{Syntax}
$\skcgen\to(pk,sk) \mathsf{Sign}(sk,m) \to \sigma, \mathsf{Vrfy}(pk,m,\sigma) \to b$
\vspace{-0.15in}
\defn{Correctness}{
    Except with negligible probability over the choice of $(pk,sk)$ output by $\skcgen(1^n)$, it must
    be that for any legal msg $m$,
    \vspace{-0.2in}
    \[\mathsf{Vrfy}(pk,m,\mathsf{Sign}(sk,m))=1\]
}
\vspace{-0.2in}
\defn{Signature EU}{
    $\mathsf{EXPT}^\mathcal{A}(1^n)$
    Challenger sends public key. $\mathcal{A}$ makes queries. Then sends its guess $(m,\sigma)$ and
    succeeds if $\mathsf{Vrfy}(pk,m,\sigma)=1$ and $m$ was not queried before.
}
\vspace{-0.2in}
\defn{Factoring}{
    $\mathsf{Factor}^\mathcal{A}(1^n)$
    Sample a random $N=pq$, where $p,q$ are $n$-bit primes.
    $\mathcal{A}$ is given $N$ and outputs two primes $p',q'$. Succeeds if $p'q'=N$.
}
\vspace{-0.2in}
\defn{RSA}{
    $\mathsf{RSA}^\mathcal{A}(1^n)$
    Sample a random $N=pq$, where $p,q$ are $n$-bit primes.
    Sample $e,d$ s.t. $ed=1 \mod \phi(N)=(p-1)(q-1)$
    Sample $y \in \Z_N^*$
    then $\mathcal{A}$ sends back an $x$ succeeds if $x^e=y \mod N$.
}
\vspace{-0.2in}
\defn{Hash, RSA, and Sign}{
    $\skcgen$ sample a random $N=pq$ s.t. $p,q$ are $n$-bit primes. Sample $e,d$ s.t. $ed=1 \mod \phi(N)=(p-1)(q-1)$.
    $sk=(N,d)$, $pk=(N,e)$ and generate a hash function.

    $\mathsf{Sign}(sk,m) = (H(m))^d \mod N = \sigma$

    $\mathsf{Vrfy}(pk,m,\sigma)$ and output 1 if $H(m)=\sigma^e \mod N$.
}
\end{document}